\handout
{5.3 Moments of Mass}
{\textbf{Chapter 5: Area and Volume Integral}}{Out: 04/02/2025}
{Author: Jaehoon Song}{Due: 04/16/2025}
{MVC: 5.3 Moments of Mass}



% \handout
% {Unit 9-1: Traditional Approach}
% {\textbf{Unit 9: Differential Equations}}{Out: 02/22/2025}
% {AP Calculus AB/BC}{Due: --/--/2025}
% {9-1: Traditional Approach}


% \input{./hw_instructions.tex}

% Define custom commands for NP and coNP
\newcommand{\NP}{\textsf{NP}}
\newcommand{\coNP}{\textsf{co-NP}}
\newcommand{\SAT}{\textsf{SAT}}
\newcommand{\coSAT}{\textsf{co-SAT}}
\newcommand{\VERTEXCOVER}{\textsf{VERTEX-COVER}}
\newcommand{\DOMINATINGSET}{\textsf{DOMINATING-SET}}
\newcommand{\SPANBOUND}{\textsf{SPAN-BOUND}}
\newcommand{\HAMPATH}{\textsf{HAM-PATH}}
















% 3D Center‑of‑Mass: Mass Density and Moments of Mass

% 3D Center‑of‑Mass: Mass Density and Moments of Mass

\section*{1. Mass Density Function}
\begin{define}{Mass Density Function}
  A function $\delta:V\subset\mathbb{R}^3\to\mathbb{R}_{\ge0}$ assigning to each point $(x,y,z)\in V$ its mass per unit volume.  In differential form, an infinitesimal mass element is $dm=\delta(x,y,z)\,dV$.  Here $dV=dx\,dy\,dz$ denotes the volume element.
\end{define}

\section*{2. Moments of Mass}

\subsection*{2.1 General $n$th Moments}
\begin{define}{General $n$th Moment about Coordinate Planes}
  For any integer $n\ge1$, the $n$th moment of mass about each coordinate plane is defined by
  \[
    M^{(n)}_{yz}=\iiint_V x^n\,\delta(x,y,z)\,dV,
    \quad
    M^{(n)}_{xz}=\iiint_V y^n\,\delta(x,y,z)\,dV,
    \quad
    M^{(n)}_{xy}=\iiint_V z^n\,\delta(x,y,z)\,dV.
  \]
  Additionally, moments about coordinate axes are defined as:
  \[
    M^{(n)}_x = \iiint_V x^n\,\delta(x,y,z)\,dV,
    \quad
    M^{(n)}_y = \iiint_V y^n\,\delta(x,y,z)\,dV,
    \quad
    M^{(n)}_z = \iiint_V z^n\,\delta(x,y,z)\,dV.
  \]
  \noindent
  Here:
  \begin{itemize}
    \item $M^{(n)}_{yz}$ refers to the $n$th moment of mass about the $yz$-plane (i.e. weighted by $x^n$).
    \item $M^{(n)}_{xz}$ refers to the $n$th moment of mass about the $xz$-plane (i.e. weighted by $y^n$).
    \item $M^{(n)}_{xy}$ refers to the $n$th moment of mass about the $xy$-plane (i.e. weighted by $z^n$).
    \item $M^{(n)}_x$ refers to the $n$th moment of mass along the $x$-axis.
    \item $M^{(n)}_y$ refers to the $n$th moment of mass along the $y$-axis.
    \item $M^{(n)}_z$ refers to the $n$th moment of mass along the $z$-axis.
  \end{itemize}
\end{define}

\subsection*{2.2 First and Second Moments}

\begin{define}{First Moment (Center‑of‑Mass Calculation)}
  The first moment ($n=1$) of mass about each plane is used to find the center‑of‑mass coordinates:
  \[
    M_{yz}=\iiint_V x\,\delta(x,y,z)\,dV,
    \quad
    M_{xz}=\iiint_V y\,\delta(x,y,z)\,dV,
    \quad
    M_{xy}=\iiint_V z\,\delta(x,y,z)\,dV.
  \]
  Dividing by the total mass $M=\iiint_V\delta(x,y,z)\,dV$ gives
  \[
    X_{\mathrm{CM}}=\frac{M_{yz}}{M},
    \quad
    Y_{\mathrm{CM}}=\frac{M_{xz}}{M},
    \quad
    Z_{\mathrm{CM}}=\frac{M_{xy}}{M}.
  \]
\end{define}

\begin{define}{Second Moment (Moment of Inertia Components)}
  The second moment ($n=2$) about each plane measures the distribution of mass and appears in the moment of inertia formulas:
  \[
    I_{yz}=\iiint_V x^2\,\delta(x,y,z)\,dV,
    \quad
    I_{xz}=\iiint_V y^2\,\delta(x,y,z)\,dV,
    \quad
    I_{xy}=\iiint_V z^2\,\delta(x,y,z)\,dV.
  \]
  In rigid‑body mechanics, the moment of inertia about the $x$‑axis is
  \[
    I_x = \iiint_V (y^2+z^2)\,\delta(x,y,z)\,dV.
  \]
\end{define}


\newpage



% 3D Moments of Mass

\section*{Moments of Mass in 3D}

\subsection*{1. Mass Density and Total Mass}

\begin{define}{Mass Density Function}
  A function $\delta:V\subset\mathbb{R}^3\to\mathbb{R}_{\ge0}$ assigning to each point $(x,y,z)\in V$ its mass per unit volume.
\end{define}

\begin{define}{Total Mass}
  The total mass of the body over volume $V$:
  \[
    M = \iiint_V \delta(x,y,z)\,\mathrm{d}V.
  \]
\end{define}

\subsection*{2. First Moment about the \emph{yz}–Plane}

\begin{define}{First Moment about the \emph{yz}–Plane}
  The first moment of mass about the $yz$–plane:
  \[
    M_{yz} = \iiint_V x\,\delta(x,y,z)\,\mathrm{d}V.
  \]
\end{define}

\subsection*{3. General $n$th Moment about the \emph{yz}–Plane}

\begin{define}{General $n$th Moment about the \emph{yz}–Plane}
  For any integer $n\ge1$, the $n$th moment of mass about the $yz$–plane is
  \[
    M^{(n)}_{yz}
    = \iiint_V x^n\,\delta(x,y,z)\,\mathrm{d}V,
  \]
  where $\mathrm{d}V = \mathrm{d}x\,\mathrm{d}y\,\mathrm{d}z$.
\end{define}







\newpage

% Moments of Mass in 3D

\section*{Moments of Mass}

\subsection*{Total Mass}
\[
  M = \iiint_V \delta(x,y,z)\,\mathrm{d}V
\]

\subsection*{First Moment about the \emph{yz}–Plane}
\[
  M_{yz} = \iiint_V x\,\delta(x,y,z)\,\mathrm{d}V
\]
\textit{(Used to compute }$X_{\mathrm{CM}} = M_{yz}/M$\textit{.)}

\subsection*{Second Moment about the \emph{yz}–Plane}
\[
  I_{yz} = \iiint_V x^2\,\delta(x,y,z)\,\mathrm{d}V
\]
\textit{(A component of the moment of inertia about the $x$–axis.)}

\subsection*{Moment of Inertia about the $x$–Axis}
\[
  I_x 
  = \iiint_V \bigl(y^2 + z^2\bigr)\,\delta(x,y,z)\,\mathrm{d}V
  \quad\bigl(\text{i.e.\ second moment about the $x$–axis / $yz$–plane}\bigr)
\]

\subsection*{General $n$th Moment about the \emph{yz}–Plane}
\[
  M^{(n)}_{yz} 
  = \iiint_V x^n\,\delta(x,y,z)\,\mathrm{d}V
\]

\newpage

% 3D Center‐of‐Mass Summary

\section*{3D Center‐of‐Mass Formulation}

\subsection*{1. Mass density and total mass}
\begin{itemize}
  \item \textbf{Density function:} $\delta(x,y,z)$ gives the mass per unit volume at each point $(x,y,z)$.
  \item \textbf{Total mass:}
    \[
      M = \iiint_V \delta(x,y,z)\,\mathrm{d}x\,\mathrm{d}y\,\mathrm{d}z.
    \]
\end{itemize}

\subsection*{2. First moments of mass}
To find the “balance point” along each axis, compute three first moments—each is a weighted integral of the distance from the corresponding coordinate plane:

\begin{itemize}
    \item \textbf{Moment $M_{yz}$:} 
        \[
            M_{yz} = \iiint_V x\,\delta(x,y,z)\,\mathrm{d}x\,\mathrm{d}y\,\mathrm{d}z
        \]
        Represents the $x$–distance from the \emph{yz}–plane.
    \item \textbf{Moment $M_{xz}$:} 
        \[
            M_{xz} = \iiint_V y\,\delta(x,y,z)\,\mathrm{d}x\,\mathrm{d}y\,\mathrm{d}z
        \]
        Represents the $y$–distance from the \emph{xz}–plane.
    \item \textbf{Moment $M_{xy}$:} 
        \[
            M_{xy} = \iiint_V z\,\delta(x,y,z)\,\mathrm{d}x\,\mathrm{d}y\,\mathrm{d}z
        \]
        Represents the $z$–distance from the \emph{xy}–plane.
\end{itemize}

\subsection*{3. Center‐of‐mass coordinates}
Each coordinate of the center of mass is the corresponding first moment divided by the total mass:
\[
  X_{\mathrm{CM}} = \frac{M_{yz}}{M}, 
  \quad
  Y_{\mathrm{CM}} = \frac{M_{xz}}{M}, 
  \quad
  Z_{\mathrm{CM}} = \frac{M_{xy}}{M}.
\]

\subsection*{Intuition}
The first moment $M_{yz} = \iiint_V x\,\delta\,\mathrm{d}V$ “weights” each element of mass by how far it sits from the $yz$–plane; dividing by $M$ gives the average $x$–position.  Analogous reasoning applies for $Y_{\mathrm{CM}}$ and $Z_{\mathrm{CM}}$.



\newpage



\begin{enumerate}
    \item \textbf{co-NP} (4 points)\\
    Show that if the complement of \textsf{CLIQUE} is in \NP, then \(\NP = \coNP\). 
    (Recall that to prove two sets are equal, show that each set is a subset of the other.)
    % % Definition of NP
    % \begin{define}{Definition: $\NP$}
    % A language $L\subseteq\{0,1\}^*$ belongs to the class $\NP$ if there exists a polynomial 
    % $p:\mathbb{N} \to \mathbb{N}$ and a deterministic Turing machine (called the \emph{verifier}) $V$ such that for every 
    % $x\in\{0,1\}^*$,
    % \[
    % x\in L \leftrightarrow \exists\, y\in\{0,1\}^{p(|x|)} \quad \mbox{with } V(x,y)=1.
    % \]
    % Here, $y$ is called a certificate (or witness) for the membership of $x$ in $L$.
    % \end{define}
    % Definition of coNP
    \begin{define}{Definition: $\coNP$}
    A language $L\subseteq\{0,1\}^*$ is in $\coNP$ if its complement
    \[
    \overline{L} = \{ x\in\{0,1\}^* : x\notin L \} \in \NP.
    \]
    In other words, for $L \in \coNP$, every "no'' instance has a short certificate that can be verified in polynomial time.
    \end{define}
    \begin{prove}
        Assuming \(\overline{\textsf{CLIQUE}}\in\NP\), by 
        the definition of \(\coNP\) it follows that \(\textsf{CLIQUE}\in\coNP\). 
        Since \textsf{CLIQUE} is \(\NP\)\nobreakdash-complete, every language 
        \(L\in\NP\) has a polynomial‑time many‑one reduction to \textsf{CLIQUE}. 
        Let $f, g \colon \{0,1\}^* \to \{\text{encodings of graphs and integer }k\}$ be
        the reductions for the following proof.
        \begin{enumerate}
            \item \(\NP \subseteq \coNP\):\\
            Let \(L\in\NP\) be arbitrary.  Then there exists a polynomial‑time reduction \(f\) such that
            \[
                x\in L \leftrightarrow f(x)\in \textsf{CLIQUE}.
            \]
            By the assumption, \(\overline{\textsf{CLIQUE}}\in\NP\), 
            \[
                x \in \overline{L} 
                \leftrightarrow
                x\notin L
                \leftrightarrow
                f(x)\notin \textsf{CLIQUE}
                \leftrightarrow
                f(x)\in \overline{\textsf{CLIQUE}},
            \]
            and so \(\overline{L}\in\NP\).  By definition of \(\coNP\), it follows that \(L\in\coNP\).  Thus \(\NP\subseteq\coNP\).
    
            \item \(\coNP \subseteq \NP\):\\
            Let \(L\in\coNP\).  Then \(\overline{L}\in\NP\) by definition.  Since \textsf{CLIQUE} is \(\NP\)\nobreakdash-complete, there exists a polynomial‑time reduction \(g\) such that
            \[
                x\in \overline{L}
                \leftrightarrow
                g(x)\in \textsf{CLIQUE}.
            \]
            By the assumption, \(\overline{\textsf{CLIQUE}}\in\NP\),  
            there is a nondeterministic polynomial‑time machine 
            \(M_{\overline{\textsf{CLIQUE}}}\) deciding \(\overline{\textsf{CLIQUE}}\). Then,
            A nondeterministic polynomial‑time machine \(M_L\) for \(L\) is constructed as follows.
            \begin{enumerate}
                \item On input \(x\), compute \(g(x)\) deterministically in polynomial time.
                \item Simulate \(M_{\overline{\textsf{CLIQUE}}}\) on \(g(x)\).
                \item Accept if and only if \(M_{\overline{\textsf{CLIQUE}}}\) accepts.
            \end{enumerate}
            For all \(x\),
            \[
                x\in L
                \leftrightarrow
                x\notin \overline{L}
                \leftrightarrow
                g(x)\notin \textsf{CLIQUE}
                \leftrightarrow
                g(x)\in \overline{\textsf{CLIQUE}}
                \leftrightarrow
                M_{\overline{\textsf{CLIQUE}}}\text{ accepts }g(x)
                \leftrightarrow
                M_L\text{ accepts }x
            \]
            Therefore, \(M_L\) decides \(L\) in \(\NP\), proving \(L \in \NP\).
            This shows \(\coNP \subseteq \NP\).
        \end{enumerate}
        Both \(\NP \subseteq \coNP\) and \(\coNP \subseteq \NP\) have been 
        proven. Thus, \(\NP = \coNP\).        
    \end{prove}
    


    % % Jaehoon
    % \item \textbf{co-NP} (4 points)\\
    % Show that if SAT is in co-NP, then NP = co-NP. 
    % (Recall that to prove two sets are equal, show that each set is 
    % a subset of the other.)
    % \begin{prove}
    %     Assuming SAT is in \coNP, by the definition of \coNP, 
    %     the complement of SAT, denoted as \(\overline{\textsf{SAT}}\), 
    %     belongs to \(\NP\). Since SAT is \(\NP\)-complete, every 
    %     language \(L \in \NP\) can be reduced to SAT in polynomial time, 
    %     that is, $L \leq_p \mathrm{SAT}$.
    %     \begin{enumerate}
    %         \item \(\NP \subseteq \coNP\): \\
    %         Let \(L \in \NP\) be any language. Since $\SAT$ is $\NP$‑complete, 
    %         there exists a polynomial‑time reduction 
    %         $f\colon \{0,1\}^* \to \{\text{CNF formulas}\}$ such that
    %         \[
    %             x\in L \leftrightarrow f(x)\in\SAT.
    %         \]
    %         Since \(\textsf{SAT} \in \coNP\), it follows that 
    %         \(\overline{\textsf{SAT}} \in \NP\). 
    %         By the same reduction \(f\) and definition of complement, 
    %         \[
    %         x \in \overline{L} \leftrightarrow x \notin L 
    %         \leftrightarrow f(x) \notin \textsf{SAT} 
    %         \leftrightarrow f(x) \in \overline{\textsf{SAT}}.
    %         \]
    %         Therefore, \(L \in \coNP\) by the definition of \coNP. 
    %         Thus, \(\NP \subseteq \coNP\).
    %         \item \(\coNP \subseteq \NP\): \\ 
    %         Let \(L \in \coNP\).  By definition of \(\coNP\), 
    %         \(\overline{L} \in \NP\).  
    %         Since \(\SAT\) is \(\NP\)‑complete, there is a 
    %         polynomial‑time reduction 
    %         \(g\colon \{0,1\}^* \to \{\text{CNF formulas}\}\) 
    %         such that
    %         \[
    %             x \in \overline{L}
    %             \leftrightarrow
    %             g(x) \in \SAT.
    %         \]
    %         Since \(\textsf{SAT} \in \coNP\), it follows that 
    %         \(\overline{\textsf{SAT}} \in \NP\). 
    %         By the same reduction \(g\) and definition of complement, 
    %         \[
    %             x \in L
    %             \leftrightarrow
    %             x \notin \overline{L}
    %             \leftrightarrow
    %             g(x) \notin \SAT
    %             \leftrightarrow
    %             g(x) \in \overline{\SAT}.
    %         \]
    %         By the definition of 
    %         \(\overline{\textsf{SAT}} \in \NP\),
    %         there exists a nondeterministic 
    %         polynomial-time machine 
    %         \(M_{\overline{\textsf{SAT}}}\) that decides 
    %         \(\overline{\textsf{SAT}}\) 
    %         constructed to decide \(L\) in polynomial 
    %         time as follows.
    %         \begin{enumerate}
    %             \item On input \(x\), \(g(x)\) is computed deterministically in polynomial time.
    %             \item \(M_{\overline{\textsf{SAT}}}\) is simulated on input \(g(x)\). 
    %             \item Accept if and only if \(M_{\overline{\textsf{SAT}}}\) accepts \(g(x)\).
    %         \end{enumerate}
    %         This machine \(M_L\) runs in nondeterministic polynomial time and for all \(x\),
    %         \[
    %             x \in L \leftrightarrow g(x) \in \overline{\textsf{SAT}} 
    %             \leftrightarrow M_{\overline{\textsf{SAT}}} \text{ accepts } g(x)
    %             \leftrightarrow M_L \text{ accepts } x.
    %         \]
    %         Therefore, \(M_L\) decides \(L\) in \(\NP\), proving \(L \in \NP\).
    %         This shows \(\coNP \subseteq \NP\).
    %     \end{enumerate}
    %     Both \(\NP \subseteq \coNP\) and \(\coNP \subseteq \NP\) have been 
    %     proven. Thus, \(\NP = \coNP\).        
    % \end{prove}






% % Definition of NP-completeness and SAT
% \begin{define}{Definition: NP-Completeness and \textsf{SAT}}
% A language $L\subseteq\{0,1\}^*$ is \emph{$\NP$-complete} if
% \begin{enumerate}
%     \item $L \in \NP$, and
%     \item For every language $L'\in\NP$, there exists a polynomial-time computable function $f$ such that 
%     \[
%     x\in L' \leftrightarrow f(x)\in L,
%     \]
%     for all $x\in \{0,1\}^*$.
% \end{enumerate}
% An example is the satisfiability problem (\textsf{SAT}), which is defined as
% \[
% \textsf{SAT}=\left\{\phi\ \big|\ \phi \text{ is a Boolean formula and } \exists\, \tau:\text{Var}(\phi) \to \{0,1\}\text{ such that } \phi(\tau)=\text{true} \right\}.
% \]
% \end{define}













    \item \textbf{Polynomial Reductions} (4 points)\\
    Show that if P = NP, then every language in P is NP-complete 
    except for $\emptyset$ and $\Sigma^*$.
    \begin{define}{Polynomial-Time Many-One Reduction}
    A language \(A\) is said to be polynomial-time many-one reducible to a language \(L\), denoted as \(A \le^p_m L\), if there exists a function \(f \in \mathrm{FP}\) such that for all \(x \in \Sigma^*\),
    \[
    x \in A \leftrightarrow f(x) \in L,
    \]
    where \(\mathrm{FP}\) denotes the class of functions computable in polynomial time.
    \end{define}
    \begin{prove}
    By direct proof, it is assumed that 
    \(\mathrm{P} = \mathrm{NP}\), and let 
    \(L \in \mathrm{P}\) such that \(L \neq \emptyset\) and 
    \(L \neq \Sigma^*\).
    It follows that \(L \in \mathrm{NP}\) by the assumption. 
    Then, the there exists 
    \(a \in L\) and \(b \in \Sigma^* \setminus L\) since 
    $\emptyset$ and $\Sigma^*$ are 
    excluded from NP-completeness. Let \(A \in \mathrm{NP}\) 
    be an arbitrary language.
    Since \(\mathrm{P} = \mathrm{NP}\), there exists a 
    deterministic polynomial-time decider \(M_A\) for \(A\) 
    with a function defined as follows. 
    \[
    f : \Sigma^* \to \Sigma^*, 
    \quad
    f(x) =
    \begin{cases}
    a, & \text{if } M_A(x) \text{ accepts},\\
    b, & \text{if } M_A(x) \text{ rejects}.
    \end{cases}
    \]
    The decider \(M_A\) runs in polynomial time, and the function 
    \(f\) outputs one of the constant-length strings \(a\) or \(b\). 
    Therefore, \(f\) belongs to the class of functions computable in polynomial time.
    \begin{enumerate}
        \item \textbf{\(x \in A \to f(x) \in L\)}: \\
        Assume \(x \in A\). By the definition of \(A\), the decider 
        \(M_A\) accepts \(x\). From the construction of \(f\), if 
        \(M_A(x)\) accepts, then \(f(x) = a\). Since \(a \in L\) 
        (by assumption), it follows that \(f(x) \in L\).
        \item \textbf{\(x \notin A \to f(x) \notin L\)}: \\
        Assume \(x \notin A\).  By definition of \(A\), the decider 
        \(M_A\) rejects \(x\).  From the construction of \(f\), 
        if 
        \(M_A(x)\) rejects, then \(f(x)=b\).  Since \(b \notin L\) (by assumption), 
        it follows that \(f(x)\notin L\).


          
    \end{enumerate}
    Therefore, it has been shown that $x \in A \leftrightarrow f(x) \in L$.

    Thus, \(f\) is a valid polynomial-time reduction from \(A\) to \(L\), that is, 
    \(A \le^p_m L\). Since \(A\) was chosen arbitrarily in \(\mathrm{NP}\), it 
    follows that \(L\) is NP-hard. (\(A\) can be arbitrarily set to a NP-hard problem.) 
    Therefore, \(L\) is NP-complete.
    Thus, under the assumption 
    \(\mathrm{P} = \mathrm{NP}\), every \(L \in \mathrm{P}\) 
    is NP‐complete except for $\emptyset$ and $\Sigma^*$.
    \end{prove}
        
        


















    \item \textbf{NP-Complete} (8 points)
    Show that the following problems are NP-complete:
    \begin{itemize}
        \item $\DOMINATINGSET$: A subset of the nodes of a graph 
        $G$ is a \textbf{dominating set} if every other node of 
        $G$ is adjacent to some node in the subset. Let 
        $$\DOMINATINGSET = \{\langle G,k \rangle | G 
        \text{ has a dominating set with } k \text{ nodes}\}$$
        Show that it is NP-complete by the reduction from $\VERTEXCOVER$.
        \begin{subprove}
        To prove that $\DOMINATINGSET$ is NP-complete, 
        it must be shown that it is in NP and it is NP-hard.\\

        \textbf{NP-membership (\(\DOMINATINGSET\in\mathrm{NP}\))}:\\
        Given an instance of \(\langle G=(V,E),k\rangle\), let \(S\subseteq V\) be a certificate 
        such that \(|S| = k\) and define a verifier running 
        in polynomial time as follows.
        \begin{enumerate}
            \item \(\lvert S\rvert = k\) is verified by counting elements of 
            \(S\) in \(O(|V|)\) time.
            \item For each \(v\in V\setminus S\), the adjacency list of \(v\) 
            is scanned in total \(O(|E|)\) time to check for the existence of 
            \(u\in S\) with \(\{u,v\}\in E\).
        \end{enumerate}
        Acceptance is granted if and only if every \(v\notin S\) 
        has a neighbor in \(S\), 
        showing a deterministic verifier of runtime \(O(|V|+|E|)\).
        Therefore, \(\DOMINATINGSET\in\mathrm{NP}\). \\

        \textbf{Reduction $f$ ($\VERTEXCOVER$ $\leq_p$ $\DOMINATINGSET$)}: \\
        From an instance 
        \(\langle G=(V,E),k\rangle\) of \(\VERTEXCOVER\) which is known to be NP‐complete, 
        Define
        \[
            G'=(V',E'),\quad k'=k,
        \]
        where
        \[
            V' = V \cup \{x_e : e\in E\},\quad
            E' = \{\{x_e,u\},\{x_e,v\}\mid e=\{u,v\}\in E\}.
        \]
        Each original edge \(e=\{u,v\}\) gives rise to a new 
        gadget‑vertex \(x_e\) adjacent only to \(u\) and \(v\).
        The transformation creates one new vertex and two new 
        edges per edge of \(G\), so 
        \(\lvert V'\rvert = \lvert V\rvert + \lvert E\rvert\) 
        and \(\lvert E'\rvert = 2\lvert E\rvert\), and runs in 
        \(O(\lvert V\rvert + \lvert E\rvert)\) time.
        \begin{enumerate}
            \item \textbf{$x\in \VERTEXCOVER \rightarrow f(x)\in\DOMINATINGSET$}\\
            Assume \(C\subseteq V\) is a vertex cover of size \(k\) in \(G\). 
            Every original vertex \(v\in V\setminus C\) is incident to some 
            edge whose other endpoint lies in \(C\), so \(v\) has a neighbor 
            in \(C\).  Each gadget‑vertex \(x_e\) for \(e=\{u,v\}\) is adjacent 
            to at least one endpoint in \(C\).  Therefore \(C\) dominates all 
            vertices of \(G'\) and is a dominating set of size \(k\).
            \item \textbf{$x\notin \VERTEXCOVER \rightarrow f(x)\notin\DOMINATINGSET$}\\
            Assume that no vertex cover of size 
            \(k\) exists in \(G\) but \(D\subseteq V'\) is a dominating set of 
            size \(k\) in \(G'\)
            for the sake of contradiction. Replace each gadget‑vertex \(x_e\in D\) by 
            one of its neighbors in \(V\), obtaining \(D'\subseteq V\) with 
            \(\lvert D'\rvert\le k\) that still dominates \(G'\).  Then every 
            gadget‑vertex \(x_e\) is dominated by an endpoint of \(e\) 
            in \(D'\), so \(D'\) covers all edges of \(G\), providing a 
            vertex cover of size \(\le k\).  
            This contradicts the assumption.  Therefore, no dominating 
            set of size \(k\) exists in \(G'\).
        \end{enumerate}
        Since $x \in$ $\VERTEXCOVER$ $\leftrightarrow f(x) \in$ $\DOMINATINGSET$ 
        in polynomial time by the reduction, 
        $\VERTEXCOVER$ $\leq_p$ $\DOMINATINGSET$. 
        Therfore, \(\DOMINATINGSET\in\mathrm{NP}\) and 
        $\VERTEXCOVER$ $\leq_p$ $\DOMINATINGSET$.
        Thus, $\DOMINATINGSET$ is NP-hard.
        \end{subprove}



        % \item $\DOMINATINGSET$: A subset of the nodes of a graph 
        % $G$ is a \textbf{dominating set} if every other node of 
        % $G$ is adjacent to some node in the subset. Let 
        % $$\DOMINATINGSET = \{\langle G,k \rangle | G 
        % \text{ has a dominating set with } k \text{ nodes}\}$$
        % Show that it is NP-complete by the reduction from $\VERTEXCOVER$.
        % \begin{subprove}
        % To prove that $\DOMINATINGSET$ is NP-complete, 
        % it must be shown that it is in NP and it is NP-hard.\\
        % \textbf{\(\DOMINATINGSET\in\mathrm{NP}\)}:\\
        % Given an instance of \(\langle G=(V,E),k\rangle\), let \(S\subseteq V\) be a certificate 
        % such that \(|S| = k\) and define a verifier running 
        % in polynomial time as follows.
        % \begin{enumerate}
        %     \item \(\lvert S\rvert = k\) is verified by counting elements of 
        %     \(S\) in \(O(|V|)\) time.
        %     \item For each \(v\in V\setminus S\), the adjacency list of \(v\) 
        %     is scanned in total \(O(|E|)\) time to check for the existence of 
        %     \(u\in S\) with \(\{u,v\}\in E\).
        % \end{enumerate}
        % Acceptance is granted iff every \(v\notin S\) has a neighbor in \(S\), 
        % showing a deterministic verifier of runtime \(O(|V|+|E|)\).
        % Therefore, \(\DOMINATINGSET\in\mathrm{NP}\). \\
        % \textbf{Reduction $f$}: \\
        % From an instance 
        % \(\langle G=(V,E),k\rangle\) of \(\VERTEXCOVER\), 
        % Define
        % \[
        %     G'=(V',E'),\quad k'=k,
        % \]
        % where
        % \[
        %     V' = V \cup \{x_e : e\in E\},\quad
        %     E' = \{\{x_e,u\},\{x_e,v\}\mid e=\{u,v\}\in E\}.
        % \]
        % Each original edge \(e=\{u,v\}\) gives rise to a new 
        % gadget‑vertex \(x_e\) adjacent only to \(u\) and \(v\).
        % The transformation creates one new vertex and two new 
        % edges per edge of \(G\), so 
        % \(\lvert V'\rvert = \lvert V\rvert + \lvert E\rvert\) 
        % and \(\lvert E'\rvert = 2\lvert E\rvert\), and runs in 
        % \(O(\lvert V\rvert + \lvert E\rvert)\) time.
        % \begin{enumerate}
        %     \item \textbf{$x \in \VERTEXCOVER \to f(x) \in \DOMINATINGSET$}\\ 
        %     If \(C\subseteq V\) is a vertex cover of size \(k\) in \(G\), 
        %     then \(C\) dominates \(G'\) because every vertex in 
        %     \(V\setminus C\) has a neighbor in \(C\), and each 
        %     gadget‑vertex \(x_e\) is adjacent to an endpoint in \(C\).
        %     \item \textbf{$x \notin \VERTEXCOVER \to f(x) \notin \DOMINATINGSET$}\\ 
        %     If \(D\subseteq V'\) is a dominating set of size \(k\) in \(G'\), any 
        %     gadget‑vertex \(x_e\in D\) may be replaced by one of its neighbors 
        %     without increasing size or losing domination, yielding 
        %     \(D'\subseteq V\) with \(\lvert D'\rvert \le k\). Since 
        %     each gadget \(x_e\) must then be dominated by an endpoint of 
        %     \(e\) in \(D'\), \(D'\) covers every edge of \(G\). 
        %     If \(\lvert D'\rvert<k\), arbitrary vertices from \(V\) 
        %     can be added to reach size \(k\).
        % \end{enumerate}
        % Since $x \in$ $\VERTEXCOVER$ $\leftrightarrow f(x) \in$ $\DOMINATINGSET$ in polynomial time, by the reduction, 
        % $\VERTEXCOVER$ $\leq_p$ $\DOMINATINGSET$. 
        % Therfore, \(\DOMINATINGSET\in\mathrm{NP}\) and $\VERTEXCOVER$ $\leq_p$ $\DOMINATINGSET$.
        % Thus, $\DOMINATINGSET$ is NP-hard.
        % \end{subprove}


            
            
        \item $\SPANBOUND$: Consider the following variant of 
        the spanning tree problem. You are given an undirected 
        graph $G = (V,E)$ and a parameter $k \in \mathbb{N}$ as 
        input. The output is YES, if there is a spanning tree 
        where every vertex of the degree has at most $k$, and 
        NO if no such spanning tree exists.
        \begin{prove}
        To prove that \(\SPANBOUND\) is NP‐complete, it must be 
        shown that it lies in NP and that it is NP‐hard.\\

        \textbf{NP-membership (\(\SPANBOUND\in\mathrm{NP}\))}:\\
        Given an instance \(\langle G=(V,E),k\rangle\), let 
        \[T\subseteq E\]
        be a certificate purported to form a spanning tree of 
        \(G\) with maximum degree at most \(k\).  A verifier 
        can check this in polynomial time as follows.
        \begin{enumerate}
            \item It is verified that \(\lvert T\rvert = \lvert V\rvert - 1\) in \(O(1)\) time.
            \item It is verified that \((V,T)\) is acyclic and connected by performing a DFS or BFS in \(O(|V|+|E|)\) time.
            \item For each \(v\in V\), the degree \(\deg_T(v)\) is counted in \(O(|V|+|E|)\) time and checked to satisfy \(\deg_T(v)\le k\).
        \end{enumerate}
        Acceptance is granted if and only if all tests succeed, 
        showing a deterministic verifier running in 
        \(O(|V|+|E|)\). Therefore, \(\SPANBOUND\in\mathrm{NP}\).\\

        \textbf{Reduction \(f\colon \HAMPATH\le_p\SPANBOUND\)}:\\
        Let 
        \[
            \HAMPATH  = \{\langle G=(V,E)\rangle\mid G\text{ has a Hamiltonian path}\},
        \]
        which is known to be NP‐complete. By an instance \(\langle G\rangle \in \HAMPATH\), define
        \[
            G' = G,
            \quad
            k' = 2.
        \]
        Clearly the transformation runs in \(O(|V|+|E|)\) time, and the output instance is \(\langle G',k'\rangle\).  It remains to show
        \[
            G\in\HAMPATH
             \leftrightarrow 
            \langle G',k'\rangle\in\SPANBOUND.
        \]
        \begin{enumerate}
            \item \(\boldsymbol{G\in\HAMPATH \rightarrow \langle G',k'\rangle\in\SPANBOUND}\).\\
            Assume that \(P\) is a Hamiltonian path in \(G\).  Then 
            \(P\) uses exactly \(\lvert V\rvert-1\) edges, spans 
            all vertices, and has maximum degree \(2\). \(P\) 
            is a spanning tree of \(G'\) since \(G'=G\) with every 
            vertex of degree at most \(2=k'\), so \(\langle G',k'\rangle\) 
            is a YES‐instance of \(\SPANBOUND\).
            \item \(\boldsymbol{\langle G',k'\rangle\in\SPANBOUND \rightarrow G\in\HAMPATH}\).\\
            Assume \(T\subseteq E\) is a spanning tree of \(G'\) 
            with maximum degree \(\le2\).  Since \(T\) is a tree 
            on \(\lvert V\rvert\) vertices with \(\lvert V\rvert-1\) 
            edges and no vertex of degree exceeding \(2\), it 
            must be a simple path that visits every vertex 
            exactly once. Therefore, \(T\) is a Hamiltonian path 
            in \(G=G'\), and so \(G\in\HAMPATH\).
        \end{enumerate}
        Since $x \in$ $\HAMPATH$ $\leftrightarrow f(x) \in$ $\SPANBOUND$ 
        in polynomial time by the reduction, 
        $\HAMPATH$ $\leq_p$ $\SPANBOUND$. 
        Therfore, \(\SPANBOUND\in\mathrm{NP}\) and 
        $\HAMPATH$ $\leq_p$ $\SPANBOUND$.
        Thus, $\SPANBOUND$ is NP-hard.
        \end{prove}
    \end{itemize}













    \item \textbf{Bonus} (1 points)
    In your opinion, what are the three most important 
    concepts or arguments or ideas presented in this course? 
    Give a 1-2 sentence explanation for why you chose each item.
    \begin{answer}
        In my opinion, the three most important concepts presented in this course in the order of 
        the right flow are like the followings.
        \begin{enumerate}
            \item \textbf{Context‑Free Languages and Turing Machines.} 
            Context‑Free Languages define exactly the patterns a pushdown 
            automaton (a parser with a stack) can recognize, showing how nested, 
            hierarchical structures are handled in computation. Building on this, 
            Turing Machines—with their infinite tape and read/write head—generalize 
            beyond any fixed memory model to capture the full scope of what an algorithm 
            can compute, setting the stage for discussions of decidability.  
            \item \textbf{Countability, Recognizability, and Undecidability.} By 
            showing that the set of all languages is uncountable while the set of 
            Turing Machines is only countable, we see that some languages cannot even 
            be listed by any algorithm (undecidable). Moreover, the distinction between 
            languages that a Turing Machine can recognize (but not necessarily decide) 
            versus those it cannot even recognize highlights the precise limits of 
            algorithmic power.  
            \item \textbf{Many‑one (\(\le_m\)) and Polynomial‑time (\(\le_p\)) Reductions.} 
            With the framework of decidability in place, many‑one reductions let us 
            translate one decision problem into another by a single computable 
            function, preserving solvability and undecidability results. Refining 
            these transformations to run in polynomial time gives us \(\le_p\) 
            reductions, which link the tractability of problems and reveal the central 
            role of NP‑complete problems in complexity theory.  
        \end{enumerate}
        
    \end{answer}
\end{enumerate}


\(A\notin\mathrm{P}\)\\
\(\text{NP-hard}\le_p A\)

